%% kapa: exam=1 round=34 year=2023 period=1 form=A subject=real_estate subject_ko=부동산학원론 seq=1 number=81
%% source: Q-Net 감정평가사 기출문제 (cst00309) artlSeq=5211789 file=raw/1008/5211789/2023년도 제34회 감정평가사 1차시험 1교시.pdf
%% conversion: kordoc 4.24.51 -> tools/gen1.py (auto); review=proofread-sonnet; flags=none
\begin{question}{81}
토지의 분류 및 용어에 관한 설명으로 옳은 것은?
\begin{choices}
  \item 획지(劃地)는 하나의 필지 중 일부에 대해서도 성립한다.
  \item 건부지(建敷地)는 건축물의 부지로 이용중인 토지 또는 건축물의 부지로 이용가능한 토지를 말한다.
  \item 나지(裸地)는 택지 중 정착물이 없는 토지로서 공법상 제한이 없는 토지를 말한다.
  \item 제내지(堤內地)는 제방으로부터 하심측으로의 토지를 말한다.
  \item 일단지(一團地)는 용도상 불가분의 관계에 있는 두 필지 이상을 합병한 토지를 말한다.
\end{choices}
\end{question}
\answer{1}
